\chapter{Data Integrity \& Storage Hierarchy}\label{chap:storage}

\section{OpenZFS}
ZFS is a file system and logical volume manager that handles high-storage environments, combined with data integrity protection. It is open-source and distributed under the CDDL (Common Development and Distribution License) license. It utilizes \texttt{Copy-On-Write} mechanism so it does not overwrite existing blocks, allowing for powerful features like the following:

\begin{itemize}
	\item \textbf{Data Integrity :} It automatically detects and restores corrupted data by using checksums and mirroring/parity of drives.
	\item \textbf{Storage Pooling :} It can combine physical drive into one large 'pool' of storage that is flexible and easy to expand in the future when adding more drives.
	\item \textbf{Snapshots \& Clones :} It creates near-instantaneous, read-only "points in time" of your data. Because ZFS is \texttt{Copy-on-Write}, these snapshots take up zero additional space until the original data is modified.
\end{itemize}

\section{Pool Initialization}
The primary storage pool, \texttt{tank}, is initialized using a RAID-Z1 topology.
This provides a balance between storage capacity and fault tolerance, allowing for a single drive failure without data loss or corruption.

\begin{lstlisting}[language=bash, caption=ZFS Pool Creation]
	# Identify disks by ID
	ls -l /dev/disk/by-id/
	
	# Create the pool with shift=12 (for 4k sectors)
	zpool create -o ashift=12 tank raidz1 \
	/dev/disk/by-id/ata-ST8000... \
	/dev/disk/by-id/ata-ST8000... \
	/dev/disk/by-id/ata-ST8000...
\end{lstlisting}


\section{Datasets Hierarchy}
ZFS datasets are probably the main part of this chapter. To maximise ZFS potential, you must put your files into different \texttt{datasets}.

Here is a basic table of datasets : \\

\noindent
\begin{tabularx}{\textwidth}{|l|X|l|l|}
\hline
\textbf{Dataset} & \textbf{Purpose} & \textbf{Recordsize} & \textbf{Snapshots} \\ \hline
\texttt{tank/appdata} & Docker config files (YAML, JSON) & 128k (Default) & Daily \\ \hline
\texttt{tank/db} & Databases (Postgres, SQLite) & 16k & Hourly \\ \hline
\texttt{tank/media} & Large files (Movies, Photos) & 1M & Weekly \\ \hline
\texttt{tank/backups} & Local backups and ISOs & 128k & None \\ \hline
\end{tabularx}

\begin{itemize}
	\item \textbf{Recordsize} is the total "write" volume every time new data is written on the disk. For general files, \texttt{128K} is fine, but for databases such as \texttt{SQL}, which writes in very small increments, it is necessary to lower the \texttt{recordsize} to prevent "write-amplification". This will increase the drive's lifespan and improves speed.
	\item \textbf{Snapshots} are used a quick-access backup, simplifying the process of returning to a past point in time without needing to use external drives. It is not a main backup, as it only concerns the data itself, not the drive. However, scheduled snapshots is a safety that is enjoyable when needed.
\end{itemize}


\section{List of ZFS Properties}

Here is a short list of properties